\slajdid{02-s025}
\begin{frame}[label=02-s025]
Znači moguće je bilo koju kombinacionu mrežu realizovati, izuzimajući invertore samo pomoću NI logičkih kola. Ili generalnije ako invertor smatramo NI logičkim kolom sa jednim ulazom: Bilo koju kombinacionu mrežu možemo realizovati korišćenjem samo NI logičkih kola. Do istog rezultata smo mogli doći i preko izraza Bulove algebre i De Morganovih obrazaca:
\[
\begin{aligned}
F&=\bar CBA+C\bar BA+CB\bar A\\
&=\overline{\overline{\bar CBA+C\bar BA+CB\bar A}}
=\overline{\overline{\bar CBA}\cdot\overline{C\bar BA}\cdot\overline{CB\bar A}}
\end{aligned}
\]
\end{frame}
